% ============================================================
% Chapter 7: Results
%   Section: Discussion
% Included from chapters/07-results/results.tex
% ============================================================

\section{Discussion}
\label{sec:discussion}

\paragraph{The steps before the verdict work better than the verdict.}
The clearest results of the project so far are about finding the news,
not about verifying it. Scoping, screening and ranking a round do what
an editor needs from them: a fraction of the links, and first places
that hold few stories not worth offering, with signals that cost no
model call (Section~\ref{sec:results-rq1}). That is consistent with the
literature's view that automation is most useful where it assists a
step of the work rather than replaces the person
\cite{nakov2021assisting}. It also rests on AI labels, and the first
measurement against a person may lower it.

\paragraph{A small local verifier is not yet a fact-checker.} On
X-Fact the 3B model is barely above chance and leans to \texttt{FALSE}
whatever the evidence. Part of that is the benchmark: claims years old,
searched on today's web, with the verdict sometimes on the page; part
is the label mapping and the stratified sample; and part is the model.
The pilot cannot separate those parts, and a stronger model on the same
claims, and the same model on constructive news, are the two
measurements that would. The design choices of the verifier --- a
closed label set, forced citation, verbatim quotes, and a confidence
recomputed from the evidence --- make its failures visible and its
right answers checkable, but they do not make a weak model strong.

\paragraph{Retrieval and reasoning fail together.} The disagreement
between the two attribution rules on the pilot (28 of 29 claims) is the
most instructive single number: the model misread what it was shown,
and it was never shown what the fact-checker used. Improving either
alone would not have fixed most of those claims. That is the case for
the oracle-evidence condition of RQ2, which gives the verifier the
annotator's own evidence and so separates the two.

\paragraph{Against the state of the art.} The supervised baselines of
the X-Fact paper reach a macro-$F_1$ of 38 to 42 on its in-domain test
set; the zero-shot pipeline reaches 11 on the Spanish part of it. The
difference is large and in the expected direction, and
Section~\ref{subsec:pilot-comparison} lists why the two are not like
for like. AVeriTeC's restriction of evidence to what existed before the
claim \cite{schlichtkrull2023averitec} and Glockner and colleagues'
demonstration of leaked evidence \cite{glockner2022missing} are both
borne out by the pilot: 23\% of the dated ranked sources were from the
year of the search, the most frequent year of all.

\paragraph{What the evaluation itself contributes.} The harness, the
custom set and the experiments are designed so that every number comes
with its interval, its provenance and the claims set apart. That is
what allowed the pilot to show that one result was a measurement of
time and not of verification, and what will make the custom-set
results interpretable when they come.

\paragraph{On constructive news the limit is the reader.} The custom
set's three evidence conditions answer what the X-Fact pilot could not
separate. The pipeline's own search, on recent claims, came within
three points of the annotator's own pages; the model alone was twenty
below. So the search is no longer the first problem; the 3B model's
reading is. Its lean to \texttt{FALSE} came from how it was asked as
much as from its size: shown every source at once it marked one as
contradicting and concluded \texttt{FALSE} by absence; asked about one
source at a time it stopped doing so, and in exchange confirmed claims
with one altered figure, because the story was the same. That trade is
the one FEVER-style substitutions are built to expose, and it is why
the set's synthetic claims are reported apart from the natural ones: a
single accuracy figure would have hidden that the adopted verifier
gained 18 points on the claims the application checks and lost 47 on
the ones built to catch it. A 9B model does both, at a cost in time on a
processor that a graphics card would remove. The replay mode is what
made these comparisons possible: two searched runs of the same pipeline
differ by six to nine points with the same recall, more than most of
the differences measured, and only fixed evidence attributes a
difference to the verifier.
\pending{the discussion of RQ1, RQ3 and RQ4 on the custom set, once the
author's labels and the annotator's review are in.}
